\documentclass[11pt,a4paper]{article}
\usepackage[margin=1in]{geometry}
\usepackage{amsmath,amssymb,amsthm}
\usepackage{graphicx}
\usepackage{booktabs}
\usepackage{enumitem}
\usepackage{hyperref}
\usepackage{natbib}
\usepackage{parskip}
\usepackage{xcolor}

\hypersetup{colorlinks=true,linkcolor=blue!60!black,citecolor=blue!60!black,urlcolor=blue!60!black}

\title{Methods in Short-Horizon Alpha Research:\\
From Statistical Arbitrage to Microstructure Signals}
\author{Mengren Man}
\date{April 2026}

\begin{document}
\maketitle

\begin{abstract}
This article surveys the principal methods used in short-horizon alpha
research for systematic trading, with emphasis on cryptocurrency
perpetual futures. We cover five broad families of approach:
(1)~mean-reversion and statistical arbitrage via cointegration and
PCA residualization, (2)~microstructure signal extraction from
trade-and-quote data, (3)~information-theoretic bar construction,
(4)~signal combination and portfolio-level stacking, and
(5)~cost-aware backtesting and the economics of edge.
Each section states the core idea, the mathematical formulation,
the practical considerations, and the key references.
The presentation is grounded in concrete examples drawn from
BTCUSDT perpetual futures on Binance.
\end{abstract}

\tableofcontents
\newpage

%% ================================================================
\section{Introduction}

Alpha research is the process of identifying predictable structure in
asset returns that can be translated into a profitable trading strategy.
At short horizons---seconds to tens of minutes---the relevant structure
is primarily \emph{microstructural}: patterns in order flow, quote
dynamics, and cross-asset co-movements that arise from the mechanics of
how markets clear, rather than from macroeconomic fundamentals.

The central challenge at short horizons is that per-trade edge is
small---typically $0.1$--$5$ basis points (bps) per round-trip---and
must be weighed against transaction costs (exchange fees, spread
crossing, market impact) that can easily exceed the edge itself.
Successful short-horizon trading therefore requires not only signal
discovery but also careful cost modelling, signal stacking for
diversification, and execution optimization.

This article surveys the principal methods in the field, organized by
the stage of the research pipeline in which they appear.

%% ================================================================
\section{Mean-Reversion and Statistical Arbitrage}

\subsection{The core idea}

Mean-reversion alpha rests on the observation that certain price
processes---or, more precisely, certain \emph{functions} of
prices---tend to revert toward a long-run mean. When the process
deviates above its mean, the expected next move is downward (and vice
versa). The alpha is the ability to predict the \emph{sign} of the
next move from the current deviation.

\subsection{Pairs trading and cointegration}

The simplest setting is two co-moving assets $X_t$ and $Y_t$ whose
prices are individually non-stationary (random walks) but whose
linear combination
\begin{equation}
  s_t = Y_t - \beta X_t
\end{equation}
is stationary. This is the definition of \emph{cointegration}
\citep{engle1987cointegration}. The stationary spread $s_t$ mean-reverts,
and the standard signal is its z-score:
\begin{equation}
  z_t = \frac{s_t - \mu_s}{\sigma_s}
\end{equation}
where $\mu_s$ and $\sigma_s$ are rolling estimates of the spread's mean
and standard deviation. The trading rule is: go long the spread when
$z_t \ll 0$ (spread is cheap), go short when $z_t \gg 0$.

\paragraph{Practical considerations.}
The hedge ratio $\beta$ is not constant. In practice it is estimated
via rolling OLS, or---more elegantly---via a \textbf{Kalman filter}
\citep{hamilton1994time}, which allows $\beta_t$ to evolve as a latent
state. The Kalman filter simultaneously estimates the current spread
and the uncertainty around $\beta_t$, providing a natural measure of
regime stability.

\subsection{PCA residualization (Avellaneda--Lee)}
\label{sec:avellaneda}

Pairs trading is the two-asset special case of a more general
framework: \textbf{statistical arbitrage via factor residualization}
\citep{avellaneda2010statistical}. Given a panel of $N$ asset returns
$R_{it}$:

\begin{enumerate}[leftmargin=1.5em]
  \item Perform PCA on the standardized return matrix. The first $K$
        principal components capture the common factors (market beta,
        sector, etc.).
  \item For each asset $i$, regress its returns on the retained PCs:
        \begin{equation}
          R_{it} = \sum_{k=1}^{K} \beta_{ik}\, \mathrm{PC}_{kt} + \varepsilon_{it}
        \end{equation}
  \item The residual $\varepsilon_{it}$ is the ``idiosyncratic'' return,
        stripped of common factor exposure.
  \item If $\varepsilon_{it}$ is mean-reverting (negative lag-1
        autocorrelation), trade the z-score of its cumulative sum
        (the ``s-score'' in Avellaneda--Lee terminology).
\end{enumerate}

This construction generalizes pairs trading in two ways: it handles
$N > 2$ assets simultaneously, and it strips \emph{all} dominant factors
(not just a single co-movement) before testing for mean-reversion.

\paragraph{Application to crypto.}
On a panel of BTC/ETH/SOL USDT perpetual futures, PC1 typically explains
$\sim$85\% of joint variance and is interpretable as the common ``crypto
market'' factor. After stripping PC1, the BTC residual exhibits lag-1
autocorrelation of $\rho_1 \approx -0.11$ to $-0.13$---significantly
more negative than the raw BTC return ($\rho_1 \approx 0$)---confirming
that the idiosyncratic component is more mean-reverting than the
total return.

\subsection{Ornstein--Uhlenbeck modelling}

The continuous-time analog of a discrete mean-reverting process is the
\textbf{Ornstein--Uhlenbeck (OU) process}:
\begin{equation}
  ds_t = \kappa\,(\theta - s_t)\,dt + \sigma\,dW_t
\end{equation}
where $\kappa > 0$ is the rate of mean-reversion, $\theta$ is the
long-run mean, and $\sigma$ is the diffusion coefficient. The OU
framework provides:
\begin{itemize}[leftmargin=1.5em]
  \item A \textbf{half-life} of mean-reversion: $t_{1/2} = \ln(2)/\kappa$.
  \item An \textbf{optimal entry/exit threshold} via dynamic programming
        \citep{leung2015optimal}.
  \item A natural connection to the Kalman filter: the OU process is a
        continuous-time AR(1), and the Kalman filter is the optimal
        recursive estimator for its parameters.
\end{itemize}

%% ================================================================
\section{Microstructure Signal Extraction}

At horizons below $\sim$1 minute, the dominant source of predictability
shifts from price-level patterns to \textbf{trade-and-quote dynamics}.
The key objects are the \emph{trade tape} (a timestamped sequence of
prices, quantities, and aggressor sides) and the \emph{order book}
(the standing bid and ask queues).

\subsection{Trade-sign autocorrelation}
\label{sec:sign-acf}

The aggressor sign of trade $t$ is $\epsilon_t \in \{-1, +1\}$
(seller-initiated or buyer-initiated). In equity markets, the
autocorrelation function of $\epsilon_t$ in \emph{trade time}
typically exhibits a slow, power-law decay:
\begin{equation}
  \rho_k \sim k^{-\gamma}, \qquad \gamma \in [0.5,\, 1.0]
  \label{eq:sign-acf}
\end{equation}
This long-memory signature reflects the prevalence of
\textbf{order-splitting}: large (institutional) orders are broken into
many child orders executed over time, creating persistent signed flow
\citep{lillo2004long, bouchaud2009markets}.

\paragraph{Empirical finding on crypto perps.}
On 7 days of BTCUSDT aggTrades, we observe $\rho_1 = +0.63$ with a
fitted decay $\gamma = 0.78$---squarely within the equity-literature
range. ETHUSDT shows $\rho_1 = +0.50$, $\gamma = 0.99$. SOLUSDT, being
less liquid (271\,ms mean inter-arrival vs.\ 66\,ms for BTC), shows the
opposite pattern at lag 1: $\rho_1 = -0.16$, the \textbf{bid--ask-bounce}
signature. Two distinct microstructure regimes are cleanly separated by
liquidity.

\subsection{Order flow imbalance (OFI)}

Aggregating signed trade volume over a short window gives the
\textbf{order flow imbalance}:
\begin{equation}
  \mathrm{OFI}_t = \sum_{i \in \text{window}} \epsilon_i \cdot q_i
\end{equation}
where $q_i$ is the quantity of trade $i$. OFI is a measure of net
buying or selling pressure. Its predictive content for future returns
has been documented extensively
\citep{cont2014price, cartea2015algorithmic}.

In a mean-reversion context, OFI is used as a \textbf{contrarian}
signal: excess buying pressure ($\mathrm{OFI} \gg 0$) is expected to
reverse as the transient demand subsides.

\subsection{Price impact}
\label{sec:impact}

The \textbf{price-impact function} measures how much the mid-price
moves in response to a signed trade of size $q$:
\begin{equation}
  \Delta p \approx \lambda \cdot f(q)
\end{equation}
where $\lambda$ is Kyle's lambda \citep{kyle1985continuous} and $f(q)$
is typically concave (square-root law for large $q$):
\begin{equation}
  f(q) \approx \sigma \sqrt{\frac{q}{V}}
\end{equation}
with $V$ the average daily volume \citep{bouchaud2009markets,
almgren2005direct}.

Understanding the impact function is essential for \textbf{cost
modelling}: any alpha signal that requires market orders will lose
a portion of its edge to impact, and the loss is a function of order
size. This is why the per-turn edge ($\sim$0.1--0.25~bps for a typical
mean-reversion signal on BTC) must be compared against the \emph{total}
cost including impact, not just exchange fees.

\subsection{Trade-clock vs.\ wall-clock sampling}

A recurring theme in microstructure research is that financial time
flows unevenly. \citet{clark1973subordinated} showed that price
variance per unit of calendar time is not constant, but scales with
the number of trades (``subordinated stochastic process''). The
practical implication: signals computed on \textbf{trade-time} or
\textbf{volume-time} bars capture more of the predictable structure
than signals on fixed wall-clock bars.

\paragraph{AR(1) comparison.}
On BTCUSDT, signed-flow AR(1) $R^2$ peaks at $\sim$0.023 in 5-second
wall-clock buckets but is $\sim$0.63 at lag-1 in trade-time. The
difference is striking: aggregating into wall-clock buckets averages
bursts with quiet periods and destroys most of the signal.

\subsection{Cross-symbol lead--lag}

In multi-asset settings, one asset's order flow may predict another
asset's price movement with a short delay. The standard diagnostic is
the \textbf{cross-correlation of signed flow} at sub-second lags.

In crypto markets, the conventional belief is that BTC flow ``leads''
altcoin flow. Empirically, this is not always detectable at 10--100\,ms
resolution on the same exchange. The lead may exist at sub-10\,ms
timescales (relevant for HFT) or in different variables (price rather
than flow).

%% ================================================================
\section{Information-Theoretic Bar Construction}

\subsection{Dollar bars and volume bars}

\citet{deprado2018advances} argues that fixed-time bars (e.g., 1-minute
OHLCV) are a poor sampling scheme for financial data because they treat
all minutes as containing equal information. Instead, he proposes
\textbf{information-driven bars}:

\begin{itemize}[leftmargin=1.5em]
  \item \textbf{Volume bars}: close a bar every $V$ units of base
        currency traded.
  \item \textbf{Dollar bars}: close a bar every $\$D$ of notional
        traded.
  \item \textbf{Tick bars}: close a bar every $N$ trades.
\end{itemize}

The motivation: by equalizing the ``information content'' per bar,
these schemes produce returns that are closer to i.i.d.\ normal,
which makes statistical tests (autocorrelation, z-scores) more
reliable.

\subsection{Empirical test}

On 7 days of BTCUSDT aggTrades, \$5M dollar bars produce a BTC
residual with lag-1 autocorrelation $\rho_1 = -0.13$ (vs.\ $-0.11$ on
1-minute time bars). The signal concentrates slightly, validating the
hypothesis, but the per-turn edge remains bounded by
$|\rho_1|\,\sigma^2 \approx 0.1$~bps---a ceiling imposed by the
horizon, not the bar construction.

%% ================================================================
\section{Signal Combination and Stacking}
\label{sec:stacking}

\subsection{The fundamental law of active management}

\citet{grinold2000active} formalize the relationship between signal
quality and portfolio performance:
\begin{equation}
  \mathrm{IR} \approx \mathrm{IC} \times \sqrt{\mathrm{BR}}
  \label{eq:fundamental-law}
\end{equation}
where IR is the information ratio (annualized Sharpe of the active
return), IC is the information coefficient (correlation between signal
and realized return), and BR is the \emph{breadth} (number of
independent bets per year).

The practical implication: a mediocre signal ($\mathrm{IC} = 0.02$)
applied to 50 instruments with daily rebalancing
($\mathrm{BR} = 50 \times 252 = 12{,}600$) achieves
$\mathrm{IR} \approx 0.02 \times \sqrt{12{,}600} \approx 2.2$---a
perfectly viable strategy despite the tiny per-bet IC.

\subsection{Stacking uncorrelated signals}

Given $N$ signals with equal Sharpe $S$ and average pairwise
correlation $\bar\rho$, the portfolio Sharpe is approximately:
\begin{equation}
  S_{\text{port}} \approx \frac{\sqrt{N}}{\sqrt{1 + (N-1)\bar\rho}}
  \cdot S
  \label{eq:stacking}
\end{equation}

When $\bar\rho \approx 0$, the improvement is $\sqrt{N}$---i.e.,
15 uncorrelated signals of Sharpe~2 produce a portfolio of Sharpe~8.
When $\bar\rho$ is large ($> 0.5$), the improvement saturates quickly
and adding more signals of the same type provides diminishing returns.

\paragraph{Empirical demonstration.}
On 12 months of BTCUSDT 1-minute klines, three signal families
(composite mean-reversion, funding-rate carry, PCA residual) have
average pairwise correlation $\bar\rho = 0.08$. The equal-weight
stack achieves a realized Sharpe multiplier of $1.71\times$, close to
the theoretical $\sqrt{3} \approx 1.73\times$.

\subsection{Weighting schemes}

Common approaches to signal weighting:
\begin{enumerate}[leftmargin=1.5em]
  \item \textbf{Equal weight}: $w_i = 1/N$. Simple, robust, hard to
        beat out of sample.
  \item \textbf{Inverse volatility}: $w_i \propto 1/\sigma_i$ where
        $\sigma_i$ is the volatility of signal $i$'s PnL. A
        risk-parity analog for signals.
  \item \textbf{Markowitz (mean-variance)}: $\mathbf{w} \propto
        \Sigma^{-1}\boldsymbol{\mu}$. Theoretically optimal but
        notoriously unstable in practice due to estimation error in
        $\Sigma$ and $\boldsymbol{\mu}$.
  \item \textbf{Shrinkage / regularized}: Ledoit--Wolf shrinkage on
        $\Sigma$, ridge/LASSO on $\boldsymbol{\mu}$.
        The practical standard at most systematic funds.
\end{enumerate}

%% ================================================================
\section{Cost-Aware Backtesting and the Economics of Edge}

\subsection{The edge-per-turn framework}

The single most important metric for a short-horizon signal is
\textbf{edge per turn} (in basis points):
\begin{equation}
  \text{edge}_{\text{bps}} = \frac{\text{gross PnL}}{\text{total turnover}}
  \times 10{,}000
\end{equation}
This number is independent of bar frequency, position sizing, and
leverage---it measures the raw informational content of the signal per
unit of trading activity.

A signal is viable if and only if:
\begin{equation}
  \text{edge}_{\text{bps}} > \text{all-in cost per turn (bps)}
\end{equation}
where all-in cost includes:
\begin{itemize}[leftmargin=1.5em]
  \item Exchange fee (taker: 1.7--4.5~bps; maker: 0--0.5~bps on
        Binance USD-M futures, depending on VIP tier)
  \item Half-spread (typically 0.01--0.1~bps for BTC perps)
  \item Market impact (a function of order size; negligible for small
        orders but dominant for large ones)
\end{itemize}

\subsection{The $|\rho_1|\sigma^2$ bound}

For a pure mean-reversion signal on a single asset at a fixed bar
frequency, the expected per-bar gross PnL of a unit-sized position is
approximately:
\begin{equation}
  \mathbb{E}[\text{gross PnL per bar}]
  \approx |\rho_1| \cdot \sigma^2
\end{equation}
where $\rho_1$ is the lag-1 autocorrelation of returns and $\sigma$ is
the per-bar volatility. This bounds the achievable edge: at 1-minute
bars on BTC ($\sigma \approx 7 \times 10^{-5}$, $|\rho_1| \approx 0.1$),
the product is $\sim$0.1~bps---an order of magnitude below taker cost.

This bound motivates three paths to profitability:
\begin{enumerate}[leftmargin=1.5em]
  \item \textbf{Maker execution}: flip the cost sign (collect spread
        instead of paying it).
  \item \textbf{Shorter horizon}: push to sub-second bars where
        $|\rho_1|$ is larger (e.g., sign ACF $\rho_1 = +0.63$ in
        trade time) even though $\sigma$ is smaller---the product
        can be 10--100$\times$ the 1-minute value.
  \item \textbf{Cross-sectional breadth}: apply the same signal to
        50+ instruments, reducing per-name turnover and gaining
        $\sqrt{N}$ diversification.
\end{enumerate}

\subsection{Walk-forward validation}

A backtest on the full sample is subject to overfitting. The standard
remedy is \textbf{walk-forward validation}: the sample is divided into
rolling train/test windows (e.g., 3-month train / 1-month test), the
signal parameters are estimated on the train window, and PnL is
evaluated on the out-of-sample test window only. A strategy that is
consistently positive across all folds is more credible than one with a
high full-sample Sharpe but variable OOS performance
\citep{bailey2014probability}.

\subsection{Paper trading as pipeline validation}

Live paper trading---running the signal computation and simulated
execution against a real-time data feed---serves as a final check that:
\begin{enumerate}[leftmargin=1.5em]
  \item The backtest implementation is correct (no look-ahead bias,
        no survivorship bias, no stale-data artifacts).
  \item The live signal distribution matches the historical one
        (no regime shift since the backtest period).
  \item The execution model's assumptions (fill at mid $\pm$
        half-spread, no queue priority, no partial fills) produce
        PnL consistent with the backtest prediction.
\end{enumerate}
Paper trading does not validate profitability---it validates the
\emph{pipeline}. If live PnL diverges sharply from the backtest
prediction, the backtest is suspect regardless of its Sharpe.

%% ================================================================
\section{Advanced Methods}
\label{sec:advanced}

\subsection{Kalman filter for time-varying hedge ratios}

In pairs trading, the hedge ratio $\beta$ drifts over time. The
\textbf{Kalman filter} models this as a linear-Gaussian state-space
system:
\begin{align}
  \text{State:} \quad & \beta_t = \beta_{t-1} + w_t, \quad w_t \sim N(0, Q) \\
  \text{Observation:} \quad & Y_t = \beta_t X_t + v_t, \quad v_t \sim N(0, R)
\end{align}
The filter recursively updates $\hat\beta_t$ and its uncertainty
$P_t$ at each new observation. The residual $v_t = Y_t - \hat\beta_t
X_t$ is the spread, and its z-score (scaled by $\sqrt{P_t + R}$)
is the trading signal.

This is more adaptive than rolling OLS and naturally down-weights
periods of high parameter uncertainty, which makes the signal more
conservative during regime transitions.

\subsection{Hidden Markov Models for regime detection}

Markets alternate between regimes (trending, mean-reverting, volatile,
quiet). A \textbf{Hidden Markov Model (HMM)} with $K$ states provides
a probabilistic estimate of the current regime:
\begin{align}
  p(S_t = k \mid \text{data}) &= \text{filtered probability}
\end{align}
The HMM parameters (transition matrix, per-state return distribution)
are estimated via the Baum--Welch algorithm. The signal is then
\emph{conditioned} on the regime: trade the mean-reversion signal only
when $p(\text{MR regime}) > \text{threshold}$.

\paragraph{Practical note.}
HMMs are powerful but fragile: with $K > 3$ states they tend to
overfit, and the Baum--Welch algorithm is sensitive to initialization.
Practitioners typically use $K = 2$ (``trending'' vs.\ ``mean-reverting'')
and re-estimate parameters monthly.

\subsection{Hawkes processes for trade arrival modelling}

Trade arrivals on crypto exchanges are bursty: clusters of rapid-fire
trades are separated by quiet intervals. The \textbf{Hawkes process}
models this self-exciting behavior:
\begin{equation}
  \lambda(t) = \mu + \sum_{t_i < t} \alpha\, e^{-\beta(t - t_i)}
\end{equation}
where $\lambda(t)$ is the intensity (arrival rate), $\mu$ is the
baseline rate, and $\alpha$, $\beta$ control the magnitude and decay
of self-excitation.

Hawkes models are used for:
\begin{itemize}[leftmargin=1.5em]
  \item Detecting periods of abnormal activity (potential signals or
        risk events).
  \item Modelling the impact of one's own trades on subsequent market
        activity (``market footprint'').
  \item Calibrating execution algorithms that need to predict
        near-future trade rates.
\end{itemize}

\subsection{Machine learning approaches}

Modern alpha research increasingly uses ML for feature selection and
non-linear signal combination:
\begin{itemize}[leftmargin=1.5em]
  \item \textbf{LASSO / Elastic Net}: sparse linear combination of a
        large feature set; useful for selecting which microstructure
        features matter.
  \item \textbf{Gradient-boosted trees (XGBoost, LightGBM)}: capture
        non-linear interactions between features without explicit
        specification; risk of overfitting, mitigated by walk-forward
        validation.
  \item \textbf{Neural networks}: sequence models (LSTM, Transformer)
        for temporal dependencies; attention mechanisms for
        cross-asset structure. State of the art in some contexts,
        but data-hungry and hard to interpret.
  \item \textbf{Meta-labeling} \citep{deprado2018advances}: a
        secondary model that predicts whether a primary signal's
        trade will be profitable, effectively learning a
        regime-conditional filter.
\end{itemize}

%% ================================================================
\section{Conclusion}

Short-horizon alpha research is not a single technique but a
\emph{pipeline}: from raw tick data, through bar construction and
signal extraction, to cost-aware backtesting and portfolio-level
stacking. The methods span classical time-series econometrics
(cointegration, PCA, AR models), market microstructure theory
(Kyle's lambda, Lo--MacKinlay sign ACF, Hasbrouck's information
share), modern portfolio theory (Grinold--Kahn, risk parity), and
machine learning (LASSO, gradient boosting, sequence models).

The unifying insight is that no single signal at the short horizon
carries enough edge to overcome transaction costs alone. The business
of systematic short-horizon trading is therefore a
\emph{portfolio construction} problem: the researcher's job is to
produce signals that are individually non-zero and mutually
uncorrelated, so that $\sqrt{N}$ does the heavy lifting at the
portfolio level.

%% ================================================================
\bibliographystyle{apalike}
\begin{thebibliography}{99}

\bibitem[Almgren et al., 2005]{almgren2005direct}
Almgren, R., Thum, C., Hauptmann, E., \& Li, H. (2005).
Direct estimation of equity market impact.
\emph{Risk}, 18(7), 58--62.

\bibitem[Avellaneda \& Lee, 2010]{avellaneda2010statistical}
Avellaneda, M. \& Lee, J.-H. (2010).
Statistical arbitrage in the US equities market.
\emph{Quantitative Finance}, 10(7), 761--782.

\bibitem[Bailey et al., 2014]{bailey2014probability}
Bailey, D.H., Borwein, J.M., de Prado, M.L., \& Zhu, Q.J. (2014).
The probability of backtest overfitting.
\emph{Journal of Computational Finance}, 20(4).

\bibitem[Bouchaud et al., 2009]{bouchaud2009markets}
Bouchaud, J.-P., Farmer, J.D., \& Lillo, F. (2009).
How markets slowly digest changes in supply and demand.
In \emph{Handbook of Financial Markets: Dynamics and Evolution}.
Elsevier.

\bibitem[Cartea et al., 2015]{cartea2015algorithmic}
Cartea, A., Jaimungal, S., \& Penalva, J. (2015).
\emph{Algorithmic and High-Frequency Trading}.
Cambridge University Press.

\bibitem[Clark, 1973]{clark1973subordinated}
Clark, P.K. (1973).
A subordinated stochastic process model with finite variance
for speculative prices.
\emph{Econometrica}, 41(1), 135--155.

\bibitem[Cont et al., 2014]{cont2014price}
Cont, R., Kukanov, A., \& Stoikov, S. (2014).
The price impact of order book events.
\emph{Journal of Financial Econometrics}, 12(1), 47--88.

\bibitem[de Prado, 2018]{deprado2018advances}
de Prado, M.L. (2018).
\emph{Advances in Financial Machine Learning}.
Wiley.

\bibitem[Engle \& Granger, 1987]{engle1987cointegration}
Engle, R.F. \& Granger, C.W.J. (1987).
Co-integration and error correction: Representation, estimation,
and testing.
\emph{Econometrica}, 55(2), 251--276.

\bibitem[Grinold \& Kahn, 2000]{grinold2000active}
Grinold, R.C. \& Kahn, R.N. (2000).
\emph{Active Portfolio Management}.
McGraw-Hill, 2nd edition.

\bibitem[Hamilton, 1994]{hamilton1994time}
Hamilton, J.D. (1994).
\emph{Time Series Analysis}.
Princeton University Press.

\bibitem[Hasbrouck, 2007]{hasbrouck2007empirical}
Hasbrouck, J. (2007).
\emph{Empirical Market Microstructure}.
Oxford University Press.

\bibitem[Kyle, 1985]{kyle1985continuous}
Kyle, A.S. (1985).
Continuous auctions and insider trading.
\emph{Econometrica}, 53(6), 1315--1335.

\bibitem[Leung \& Li, 2015]{leung2015optimal}
Leung, T. \& Li, X. (2015).
Optimal mean reversion trading with transaction costs and
stop-loss exit.
\emph{International Journal of Theoretical and Applied Finance},
18(3), 1550020.

\bibitem[Lillo et al., 2004]{lillo2004long}
Lillo, F., Mike, S., \& Farmer, J.D. (2004).
Theory for long memory in supply and demand.
\emph{Physical Review E}, 71(6), 066122.

\end{thebibliography}

\end{document}
